\section{Discussion}
The proposal is feasible within one semester because it combines existing, documented technologies (ARCore, LiteRT, WebSocket) rather than developing new tracking or networking methods, and because the single-placer scope avoids the problem of aligning several phones. The main technical risk is accuracy over distance: the evidence shows error growing with the distance from the reference~\cite{chalhoub2019}, \cite{feigl2020}, and no study measured phone AR on asphalt, a low-texture surface that is likely to increase drift~\cite{scargill2021}. Outdoor lighting, glare on the screen and wet weather during winter tests may also affect both tracking and the cone check, whose training data come from car-mounted cameras~\cite{vodisch2022}. For car identification, the main risks are the size of the self-made dataset and the similarity between the team's cars; a model that relies on the background or the light of one day rather than on the car would pass a random split and fail on a new day, which is why the test set is split by day. The user's safety is a further concern, as attention is on the screen while walking; the app is meant to be used only on a closed track.

The evidence itself has limits. Most accuracy studies come from construction, indoors or over small areas, with few participants~\cite{chalhoub2019}, \cite{kim2025isarc}; one tracking study is a preprint~\cite{scargill2021}; and the protocol comparison is a single controlled test~\cite{ajayi2023}. The review changed the proposal in two ways: a single start marker was replaced by two-point alignment with re-alignment markers, and the Geospatial API was moved from a candidate to a rejected option. The twin is also limited: dashboard edits reach the phone automatically, but a person must move the physical cone. Unity, iOS and dedicated surveying instruments were rejected for this version; collaborative placement with shared anchors and an integrated RTK receiver are left for future work. The expected contribution is a justified integration and evaluation of existing technologies for a task that has not been studied, not a new tracking method.
